% Apêndices contêm material elaborado pelo próprio autor.
% Caso deseje página divisória própria, descomente \partapendices.
\begin{apendicesenv}
  % \partapendices
  \chapter{Título do apêndice}
  Inclua aqui o material produzido pelo autor que complementa o trabalho (questionários, códigos-fonte, deduções extensas, entre outros).

% --------------------------------------------------------------------------
% Exemplo opcional de inserção de códigos-fonte (pacote listings)
% Linguagens prontas para uso: Python, R, SQL, Matlab e VBA.
% Acentos UTF-8 em português são suportados diretamente nos comentários e
% strings. Caso utilize códigos-fonte no TCC, lembre-se de descomentar
% \tagstructbegin{tag=TOC}
\listofcodigos*
\tagstructend
\cleardoublepage
 em tcc-depro-ufs.tex.
% --------------------------------------------------------------------------
%
% 1) Código escrito diretamente no arquivo .tex:
%
% \begin{lstlisting}[
%   language=Python,
%   caption={Cálculo de indicador de eficiência operacional},
%   label={cod:exemplo-python}
% ]
% # Cálculo da eficiência média com comentários em UTF-8 (ação, produção)
% def calcular_eficiencia(producao_real, capacidade_nominal):
%     if capacidade_nominal <= 0:
%         raise ValueError("Capacidade deve ser maior que zero.")
%     return (producao_real / capacidade_nominal) * 100.0
%
% print(f"Eficiência: {calcular_eficiencia(420, 500):.1f}%")
% \end{lstlisting}
% \fontecodigo{Elaboração própria}
%
% 2) Importação de arquivo externo (ex.: Python, R, SQL, Matlab ou VBA):
%
% \lstinputlisting[
%   language=VBA,
%   caption={Rotina de otimização no Excel Solver via macro VBA},
%   label={cod:exemplo-vba}
% ]{dados/macro-solver.bas}
% \fontecodigo{Elaboração própria}
\end{apendicesenv}
